% GENERATED FILE. Edit canonical lessons, then run npm run generate:books.
% canonical-source-hash: fnv1a64:51e0b24e5dc0740d
% canonical-lessons: ML-C527-dayayulla, ML-C527-ksamayulla, ML-C527-pisukkan, ML-C527-udaramaya, ML-C527-visvastan

\chapter{Kind and Stingy}
\label{ch:ml-kind-and-stingy}
\hlchaptermodality{527}

\begin{chapteropening}
\textbf{By the end of this chapter:} \emph{I can say kind-hearted, patient by nature, a miser, generous and a loyal person.}

Complete the last lesson of chapter 527: I can say kind-hearted, patient by nature, a miser, generous and a loyal person.
\end{chapteropening}

\section[\texorpdfstring{\ml{ദയയുള്ള} (dayayuḷḷa)}{ദയയുള്ള (dayayuḷḷa)}]{\ml{ദയയുള്ള} (dayayuḷḷa) — kind-hearted; kind}

\label{lesson:ML-C527-dayayulla}



\begin{quote}
Before the new one: say the Malayalam for to be amused, then the Malayalam for to grow weary.
\end{quote}

\subsection*{You'll want to know: \ml{ദയയുള്ള}}
\textbf{\ml{ദയയുള്ള}} — \emph{dayayuḷḷa} — \textquotedblleft{}kind-hearted; kind\textquotedblright{}.

\begin{quote}
\ml{ദയയുള്ള} \ml{ഒരു} \ml{സ്ത്രീ} \ml{എനിക്ക്} \ml{വഴി} \ml{പറഞ്ഞുതന്നു}.
\end{quote}

\begin{quote}
\emph{dayayuḷḷa oru strī enikkŭ vaḻi paṟaññutannu.}
\end{quote}

\begin{quote}
A kind woman told me the way.
\end{quote}

\ml{ദയ}, kindness, + \ml{ഉള്ള}, having. \ml{ദയവായി}, please, holds the same noun.

\begin{cousinweb}
\textbf{\ml{ദയയുള്ള}} (\emph{dayayuḷḷa}) is built from \ml{ദയ}, from Sanskrit \emph{dayā}, \textquotedblleft{}sympathy, compassion\textquotedblright{} (CDIAL 6185), and the native ending -\ml{ഉള്ള}, \textquotedblleft{}having\textquotedblright{} (DEDR 697).
\end{cousinweb}

\begin{grammarlens}[title={What you've built}]
One more word for describing what you lived and felt.
\end{grammarlens}

\subsection*{Your turn}
\begin{itemize}
\raggedright
  \item \emph{Say it:} \emph{dayayuḷḷa}
  \item \emph{Say it:} \emph{dayayuḷḷa}, once more
  \item \emph{Say it:} \emph{dayayuḷḷa oru strī enikkŭ vaḻi paṟaññutannu.} — A kind woman told me the way.
  \item \emph{Recall:} say the Malayalam for to be amused, then the Malayalam for to grow weary, then say \emph{dayayuḷḷa} again
\end{itemize}

\begin{tcolorbox}[breakable,colback=teal!4,colframe=teal!35!black,title={Before you move on}]
What does \ml{ദയയുള്ള} mean? (\textquotedblleft{}Kind-hearted; kind\textquotedblright{}.) Say it once more.
\end{tcolorbox}
\section[\texorpdfstring{\ml{ക്ഷമയുള്ള} (kṣamayuḷḷa)}{ക്ഷമയുള്ള (kṣamayuḷḷa)}]{\ml{ക്ഷമയുള്ള} (kṣamayuḷḷa) — patient by nature}

\label{lesson:ML-C527-ksamayulla}



\begin{quote}
Before the new one: say the Malayalam for to grow weary, then the Malayalam for kind-hearted.
\end{quote}

\subsection*{You'll want to know: \ml{ക്ഷമയുള്ള}}
\textbf{\ml{ക്ഷമയുള്ള}} — \emph{kṣamayuḷḷa} — \textquotedblleft{}patient by nature\textquotedblright{}.

\begin{quote}
\ml{ക്ഷമയുള്ള} \ml{അധ്യാപകൻ} \ml{ചോദ്യം} \ml{വീണ്ടും} \ml{വിശദീകരിച്ചു}.
\end{quote}

\begin{quote}
\emph{kṣamayuḷḷa adhyāpakan cōdyaṁ vīṇṭuṁ viśadīkariccu.}
\end{quote}

\begin{quote}
The patient teacher explained the question again.
\end{quote}

\ml{ക്ഷമ}, patience, + \ml{ഉള്ള}; \ml{ക്ഷമിക്കണം}, sorry, comes from the same Sanskrit word.

\begin{cousinweb}
\textbf{\ml{ക്ഷമയുള്ള}} (\emph{kṣamayuḷḷa}) is built from \ml{ക്ഷമ}, from Sanskrit \emph{kṣamā}, \textquotedblleft{}patience, forbearance\textquotedblright{} (CDIAL 3659), and the native ending -\ml{ഉള്ള}, \textquotedblleft{}having\textquotedblright{} (DEDR 697).
\end{cousinweb}

\begin{grammarlens}[title={What you've built}]
One more word for describing what you lived and felt.
\end{grammarlens}

\subsection*{Your turn}
\begin{itemize}
\raggedright
  \item \emph{Say it:} \emph{kṣamayuḷḷa}
  \item \emph{Say it:} \emph{kṣamayuḷḷa}, once more
  \item \emph{Say it:} \emph{kṣamayuḷḷa adhyāpakan cōdyaṁ vīṇṭuṁ viśadīkariccu.} — The patient teacher explained the question again.
  \item \emph{Recall:} say the Malayalam for to grow weary, then the Malayalam for kind-hearted, then say \emph{kṣamayuḷḷa} again
\end{itemize}

\begin{tcolorbox}[breakable,colback=teal!4,colframe=teal!35!black,title={Before you move on}]
What does \ml{ക്ഷമയുള്ള} mean? (\textquotedblleft{}Patient by nature\textquotedblright{}.) Say it once more.
\end{tcolorbox}
\section[\texorpdfstring{\ml{പിശുക്കൻ} (piśukkan)}{പിശുക്കൻ (piśukkan)}]{\ml{പിശുക്കൻ} (piśukkan) — a miser; a stingy person}

\label{lesson:ML-C527-pisukkan}



\begin{quote}
Before the new one: say the Malayalam for kind-hearted, then the Malayalam for patient by nature.
\end{quote}

\subsection*{You'll want to know: \ml{പിശുക്കൻ}}
\textbf{\ml{പിശുക്കൻ}} — \emph{piśukkan} — \textquotedblleft{}a miser; a stingy person\textquotedblright{}.

\begin{quote}
\ml{ആ} \ml{പിശുക്കൻ} \ml{ആർക്കും} \ml{ഒന്നും} \ml{കൊടുക്കില്ല}.
\end{quote}

\begin{quote}
\emph{ā piśukkan ārkkuṁ onnuṁ koṭukkilla.}
\end{quote}

\begin{quote}
That miser will not give anything to anyone.
\end{quote}

A woman is \ml{പിശുക്കി}; the quality is \ml{പിശുക്ക്}, stinginess.

\begin{cousinweb}
\textbf{\ml{പിശുക്കൻ}} (\emph{piśukkan}) is native Dravidian: DEDR 4144 lists Malayalam \ml{പിശുക്ക്}, \textquotedblleft{}stinginess\textquotedblright{}, and \ml{പിശുക്കൻ}, \textquotedblleft{}a miser\textquotedblright{} (DEDR 4144).
\end{cousinweb}

\begin{grammarlens}[title={What you've built}]
One more word for describing what you lived and felt.
\end{grammarlens}

\subsection*{Your turn}
\begin{itemize}
\raggedright
  \item \emph{Say it:} \emph{piśukkan}
  \item \emph{Say it:} \emph{piśukkan}, once more
  \item \emph{Say it:} \emph{ā piśukkan ārkkuṁ onnuṁ koṭukkilla.} — That miser will not give anything to anyone.
  \item \emph{Recall:} say the Malayalam for kind-hearted, then the Malayalam for patient by nature, then say \emph{piśukkan} again
\end{itemize}

\begin{tcolorbox}[breakable,colback=teal!4,colframe=teal!35!black,title={Before you move on}]
What does \ml{പിശുക്കൻ} mean? (\textquotedblleft{}A miser; a stingy person\textquotedblright{}.) Say it once more.
\end{tcolorbox}
\section[\texorpdfstring{\ml{ഉദാരമായ} (udāramāya)}{ഉദാരമായ (udāramāya)}]{\ml{ഉദാരമായ} (udāramāya) — generous}

\label{lesson:ML-C527-udaramaya}



\begin{quote}
Before the new one: say the Malayalam for patient by nature, then the Malayalam for a miser.
\end{quote}

\subsection*{You'll want to know: \ml{ഉദാരമായ}}
\textbf{\ml{ഉദാരമായ}} — \emph{udāramāya} — \textquotedblleft{}generous\textquotedblright{}.

\begin{quote}
\ml{ഉദാരമായ} \ml{മനസ്സുള്ള} \ml{ആളാണ്} \ml{എന്റെ} \ml{അമ്മാവൻ}.
\end{quote}

\begin{quote}
\emph{udāramāya manassuḷḷa āḷāṇŭ enṟe ammāvan.}
\end{quote}

\begin{quote}
My uncle is a man with a generous heart.
\end{quote}

The noun is \ml{ഉദാരത}, generosity. -\ml{ആയ} makes an adjective of a Sanskrit noun.

\begin{cousinweb}
\textbf{\ml{ഉദാരമായ}} (\emph{udāramāya}) is built from Sanskrit \emph{udāra}, \textquotedblleft{}lofty, noble\textquotedblright{} (CDIAL 1935), with the Malayalam adjective ending -\ml{ആയ}. Being noble came to mean giving freely.
\end{cousinweb}

\begin{grammarlens}[title={What you've built}]
One more word for describing what you lived and felt.
\end{grammarlens}

\subsection*{Your turn}
\begin{itemize}
\raggedright
  \item \emph{Say it:} \emph{udāramāya}
  \item \emph{Say it:} \emph{udāramāya}, once more
  \item \emph{Say it:} \emph{udāramāya manassuḷḷa āḷāṇŭ enṟe ammāvan.} — My uncle is a man with a generous heart.
  \item \emph{Recall:} say the Malayalam for patient by nature, then the Malayalam for a miser, then say \emph{udāramāya} again
\end{itemize}

\begin{tcolorbox}[breakable,colback=teal!4,colframe=teal!35!black,title={Before you move on}]
What does \ml{ഉദാരമായ} mean? (\textquotedblleft{}Generous\textquotedblright{}.) Say it once more.
\end{tcolorbox}
\section[\texorpdfstring{\ml{വിശ്വസ്തൻ} (viśvastan)}{വിശ്വസ്തൻ (viśvastan)}]{\ml{വിശ്വസ്തൻ} (viśvastan) — a loyal person; a trusted one}

\label{lesson:ML-C527-visvastan}



\begin{quote}
Before the new one: say the Malayalam for a miser, then the Malayalam for generous.
\end{quote}

\subsection*{You'll want to know: \ml{വിശ്വസ്തൻ}}
\textbf{\ml{വിശ്വസ്തൻ}} — \emph{viśvastan} — \textquotedblleft{}a loyal person; a trusted one\textquotedblright{}.

\begin{quote}
\ml{അവൻ} \ml{മുതലാളിയുടെ} \ml{വിശ്വസ്തനാണ്}.
\end{quote}

\begin{quote}
\emph{avan mutalāḷiyuṭe viśvastanāṇŭ.}
\end{quote}

\begin{quote}
He is the boss's trusted man.
\end{quote}

A woman is \ml{വിശ്വസ്ത}; the quality is \ml{വിശ്വസ്തത}, loyalty.

\begin{cousinweb}
\textbf{\ml{വിശ്വസ്തൻ}} (\emph{viśvastan}) is borrowed from Sanskrit \emph{viśvasta}, \textquotedblleft{}trusted, faithful\textquotedblright{} (Monier-Williams), with the Malayalam masculine ending -\ml{ൻ}. It shares its root with \ml{വിശ്വാസം}.
\end{cousinweb}

\begin{grammarlens}[title={What you've built}]
One more word for describing what you lived and felt.
\end{grammarlens}

\subsection*{Your turn}
\begin{itemize}
\raggedright
  \item \emph{Say it:} \emph{viśvastan}
  \item \emph{Say it:} \emph{viśvastan}, once more
  \item \emph{Say it:} \emph{avan mutalāḷiyuṭe viśvastanāṇŭ.} — He is the boss's trusted man.
  \item \emph{Recall:} say the Malayalam for a miser, then the Malayalam for generous, then say \emph{viśvastan} again
\end{itemize}

\begin{tcolorbox}[breakable,colback=teal!4,colframe=teal!35!black,title={Before you move on}]
What does \ml{വിശ്വസ്തൻ} mean? (\textquotedblleft{}A loyal person; a trusted one\textquotedblright{}.) Say it once more.
\end{tcolorbox}
